\section{引言}

随着数字经济发展和人工智能技术扩散，企业在公开披露中越来越频繁地使用“数字化转型”“智能化升级”“大模型赋能”等表述。尤其在 IPO 阶段，招股说明书与问询回复承担着企业向市场和监管者集中展示经营能力、技术基础和成长逻辑的重要功能。在这一场景下，AI 与数字化相关表述一方面可能是企业真实能力和长期投入的信号，另一方面也可能演化为一种叙事包装，用于强化成长故事、改善市场印象或迎合政策与产业热点。

从现实观察看，不同企业对 AI 与数字化的披露存在显著差异。有的企业能够清楚说明技术应用于哪些业务场景、投入了哪些资源、建立了哪些系统，以及这些举措如何作用于生产、营销、供应链和管理流程；也有企业主要停留在战略愿景、行业趋势或宏大叙事层面，文本中充满“全面赋能”“深度融合”“行业领先”等概念化表述，但缺少可验证的应用证据。更值得关注的是，在交易所问询过程中，监管者往往会进一步追问企业相关披露的先进性、合理性、必要性与披露充分性。这使得 IPO 阶段的文本不仅是企业自我叙事的载体，也成为观察企业真实数字化能力与叙事包装之间张力的重要窗口。

既有研究已经分别讨论了企业数字化转型与绩效、企业文本披露与信息质量、IPO 招股书文本与资本市场后果等问题。数字化转型研究强调技术应用、组织变革和能力重构对企业长期竞争力的重要作用 \citep{Vial_2019,Verhoef_2021,Warner_2019}；文本披露研究则表明，企业叙事的语气、可读性、具体性和包装方式会影响信息质量与外部认知 \citep{Healy_2001,LOUGHRAN_2016,Li_2008,Huang2014,MerklDaviesConcept2011}；IPO 文本研究进一步指出，招股书语言包含重要的信息含量，可以影响投资者判断和发行过程 \citep{Hanley_2010,Loughran_2013}。然而，较少有研究把三条线索放到同一个框架下：在 IPO 阶段，企业关于 AI 与数字化的叙事究竟是真实转型信号，还是概念包装？这种差异能否在上市后的经营结果中得到验证？

围绕上述问题，本文拟以 A 股 IPO 企业为总体范围，筛选在招股说明书或问询回复中明确提及 AI 或数字化内容的企业作为核心样本。本文首先通过关键词体系识别相关文本段落，然后采用 LLM 辅助摘要与初步编码，并在此基础上进行人工复核，重点围绕提及强度、场景具体性和表述审慎性三个维度构建 AI 披露真实性综合指数，同时保留叙事包装性的反向指标。随后，本文将企业上市后第一年的 ROA、营业收入增长率和总资产周转率作为主要经营结果变量，使用基础回归模型检验披露真实性与经营表现之间的关系。

本文的研究价值主要体现在三个方面。第一，在研究场景上，本文将数字化叙事问题前移到 IPO 阶段，扩展了企业数字化披露研究的场景边界。第二，在方法上，本文尝试把关键词识别、LLM 辅助标注和人工复核结合起来，形成一套适合公开披露文本的轻量化研究路径。第三，在现实层面，本文有助于理解企业数字化叙事的可信度问题，并为投资者识别概念包装、为监管者完善问询关注点、为企业规范信息披露提供参考。

本文余下结构安排如下：第二部分回顾相关文献并提出研究假设；第三部分说明样本、数据来源、变量设计和研究方法；第四部分给出实证结果报告模板与稳健性分析安排；第五部分总结全文并讨论研究启示。
